\part{The Interface}

\chapter{Configuration de l'écran}

FairWindSK garde trois zones de coquille visibles en tout temps.
Ces noms sont utilisés systématiquement dans le présent manuel, dans le code source du projet,
et dans toute la documentation du contributeur.

\section{Les trois zones}

\begin{center}
\begin{tcolorbox}[enhanced,colback=LtBlue,colframe=Navy,boxrule=1.5pt,arc=4pt,
  top=7pt,bottom=5pt,left=12pt,right=12pt,width=0.96\textwidth]
  {\sffamily\bfseries\color{Navy}TOP BAR --- always visible}\\[3pt]
  {\small\icon{lcd-app}~App icon \quad
         \icon{layout-signalk}~Data widgets \quad
         App name \quad
         \icon{lcd-clock}~Date/time \quad
         Status icons}
\end{tcolorbox}
\vspace{2pt}
\begin{tcolorbox}[enhanced,colback=white,colframe=Navy,boxrule=1.5pt,arc=4pt,
  top=26pt,bottom=26pt,left=12pt,right=12pt,width=0.96\textwidth]
  \centering
  {\large\sffamily\bfseries\color{MidGray}APPLICATION AREA}\\[4pt]
  {\small\sffamily\color{MidGray}\itshape
    Launcher \textbar{} Chart plotter \textbar{} Instruments \textbar{}
    Settings \textbar{} My~Data \textbar{} Signal~K web apps}
\end{tcolorbox}
\vspace{2pt}
\begin{tcolorbox}[enhanced,colback=LtBlue,colframe=Navy,boxrule=1.5pt,arc=4pt,
  top=7pt,bottom=5pt,left=12pt,right=12pt,width=0.96\textwidth]
  {\sffamily\bfseries\color{Navy}BOTTOM BAR --- always visible}\\[3pt]
  {\small
    Open web apps strip \quad
    \icon{applications}~Apps \quad
    \icon{alarm-pob}~POB \quad
    \icon{command-autopilot}~Autopilot \quad
    \icon{anchor-iec}~Anchor \quad
    \icon{alarm-general}~Alarms \quad
    \icon{database}~My~Data \quad
    \icon{settings-iec}~Settings \quad
    \icon{layout-signalk}~Status}
\end{tcolorbox}
\end{center}

\screenshot{screen-layout-overview}{Mise en page de l'écran: Barre supérieure, zone d'application, barre inférieure}{}

\section{La zone de demande}

La zone d'application est le principal espace de travail entre les barres.
C'est ici que FairWindSK montre:

\begin{itemize}
  \item launcher pages
  \item chart and navigation apps
  \item Settings
  \item My~Data
  \item embedded Signal~K web applications
\end{itemize}

La vue active ne doit pas chevaucher visuellement la barre supérieure ou la barre inférieure à moins que
un tiroir formellement défini le couvre temporairement.

\section{Tiroirs}

\rowcolors{2}{LtGray}{white}
\begin{longtable}{L{4.0cm}L{9.7cm}}
\toprule
\rowcolor{Navy}
\color{white}\sffamily\bfseries Drawer &
\color{white}\sffamily\bfseries Description \\
\midrule
Tiroir horizontal de la barre inférieure&
Boîte de dialogue modale s'ouvrant vers le haut depuis la barre inférieure.
Utilisé pour les dialogues modaux web, les confirmations, les sélectionneurs de fichiers et les sous-panneaux de paramètres.
Une seule boîte de dialogue est visible à la fois; pile de dialogues supplémentaires.
Peut s'étendre jusqu'à couvrir temporairement la zone de demande.
  The hosted UI must fit inside the drawer without requiring vertical scrolling. \\
Champ d'application Tiroir vertical&
Panneau non modal ouvrant vers la gauche du bord droit de la zone d'application.
Utilisé par les applications pour les menus contextuels ou les palettes d'outils.
  The Application Area remains interactive behind it. \\
\bottomrule
\end{longtable}

\section{Modèle à guichet unique}

FairWindSK applique une conception à guichet unique sur toutes les plateformes.
Aucune application n'ouvre une deuxième fenêtre; aucune application native externe n'est lancée.
\code{file://} launcher entries are blocked.

Android est l'exception de plate-forme intentionnelle pour les applications natives sélectionnées.
Sur Android 13/API 33 ou plus récent, FairWindSK peut en option agir comme l'application
launch activities explicitly selected in \ui{Settings~> Android}. Those activities
quitter temporairement la fenêtre FairWindSK et revenir par Android Back ou Home;
Tout le contenu FairWindSK et Signal~K continue d'utiliser le shell à guichet unique.

Sur les plateformes de bureau (Linux, macOS, Windows, Raspberry~Pi~OS),
press \key{SHIFT+TAB} to bring FairWindSK back to the foreground if a hosted web
application a pris toute l'attention.

\chapter{La barre supérieure}
\label{ch:topbar}

La barre supérieure est toujours visible.
Son but : fournir une vérification de confiance de deux secondes à tout moment.
Use it to answer: \textit{What is happening right now, and is the data still
trustworthy?}

\screenshot{topbar-overview}{Barre supérieure: position, SOG, DPT, BTW, ETA et icônes d'état}{}

\section{Éléments de barre supérieure fixes}

\rowcolors{2}{LtGray}{white}
\begin{longtable}{C{1.5cm}L{2.8cm}L{9.4cm}}
\toprule
\rowcolor{Navy}
\color{white}\sffamily\bfseries Icon &
\color{white}\sffamily\bfseries Element &
\color{white}\sffamily\bfseries Function \\
\midrule
\iconlg{lcd-app}       & App icon    &
Icône de la demande actuelle.
  Tapping opens the Application Area vertical drawer. \\
---                     &Nom de l'application&
  Name of the currently active application. \\
\iconlg{lcd-clock}     & Date / time &
  Current device time and date. \\
\iconlg{layout-signalk}& Status icons &
  REST and stream health indicators. \\
\bottomrule
\end{longtable}

\section{Widgets de données --- Référence complète}

Chaque widget s'inscrit à un chemin Signal~K configuré dans
\ui{Settings~> Signal~K Paths}.

\rowcolors{2}{LtGray}{white}
\begin{longtable}{C{1.4cm}C{1.1cm}L{4.4cm}L{3.0cm}L{3.8cm}}
\toprule
\rowcolor{Navy}
\color{white}\sffamily\bfseries Icon &
\color{white}\sffamily\bfseries Key &
\color{white}\sffamily\bfseries Default Signal~K path &
\color{white}\sffamily\bfseries Displays &
\color{white}\sffamily\bfseries Operational significance \\
\midrule
\iconlg{lcd-position} & \code{pos} & \skpath{navigation.position}                              & Latitude / longitude      & GPS fix confirmation \\
\iconlg{lcd-cog}      & \code{cog} & \skpath{navigation.courseOverGroundTrue}                  & COG (\textdegree{}T)      & Actual track over seabed \\
\iconlg{lcd-sog}      & \code{sog} & \skpath{navigation.speedOverGround}                       & SOG (kn)                  & True speed for ETAs \\
\iconlg{lcd-hdg}      & \code{hdg} & \skpath{navigation.headingTrue}                           & Heading (\textdegree{}T)  & Where the bow points \\
\iconlg{lcd-stw}      & \code{stw} & \skpath{navigation.speedThroughWater}                     & STW (kn)                  & Sail trim, current deduction \\
\iconlg{lcd-dpt}      & \code{dpt} & \skpath{environment.depth.belowTransducer}                & Depth below transducer    & Safety-critical in shoal water \\
\iconlg{lcd-btw}      & \code{btw} & \skpath{navigation.course.calcValues.bearingTrue}         & Bearing to WPT            & Course to steer \\
\iconlg{lcd-dtg}      & \code{dtg} & \skpath{navigation.course.calcValues.distance}            & Distance to go            & Passage progress \\
\iconlg{lcd-ttg}      & \code{ttg} & \skpath{navigation.course.calcValues.timeToGo}            & Time to go                & ETA planning \\
\iconlg{lcd-eta}      & \code{eta} & \skpath{navigation.course.calcValues.estimatedTimeOfArrival} & Estimated arrival      & Crew notification \\
\iconlg{lcd-xte}      & \code{xte} & \skpath{navigation.course.calcValues.crossTrackError}     & Cross-track error         & Route compliance \\
\iconlg{lcd-vmg}      & \code{vmg} & \skpath{performance.velocityMadeGood}                     & Velocity made good        & True sailing efficiency \\
\iconlg{lcd-wpt}      & \code{wpt} & \skpath{navigation.course.nextPoint}                      & Active waypoint name       & Route awareness \\
\bottomrule
\end{longtable}

\begin{warningbox}
Missing depth is \textbf{not} zero depth.
Manquer signifie que la mesure n'est pas disponible --- la traiter comme totalement inconnue.
Si la profondeur devient discontinue ou manquante dans les eaux peu profondes ou restreintes, vérifier
immédiatement et ajuster la vitesse et la marge de manœuvre en conséquence.
\end{warningbox}

\section{Lecture des valeurs de navigation}

\subsection{COG et SOG}

COG est en cours de route au sol; SOG est en vitesse au-dessus du sol.
Ceux-ci montrent comment le bateau se déplace réellement sur la terre, qui peut différer de
où l'arc pointe en raison de la légion et du courant de marée.

\subsection{Rubrique}

Le cap vous indique où l'arc est pointé.
Si le cap devient caduque ou manquant, dirigez avec prudence et contre-vérifiez
Une boussole fiable.

\subsection{Valeurs de parcours et de point de cheminement}

Lorsque le routage est actif, BTW, DTG, TTG, ETA, XTE et VMG mettent à jour en permanence.
Si l'une de ces valeurs cesse d'être mise à jour, vérifiez le graphique et votre environnement avant
s'appuyant sur le guidage.

\section{Configuration de la barre supérieure}

Open \ui{Settings~> Top Bar}.
Un aperçu en direct de la barre apparaît en haut; la palette du widget est en dessous.
Use the \iconlg{arrow-left-google}~move-left and \iconlg{arrow-right-google}~move-right
boutons pour repositionner les widgets.
Toggle \key{Show Icon}, \key{Show Text}, \key{Show Units}, and \key{Show Trend}
par widget.
Tap \iconlg{refresh-google}~\key{Reset Defaults} to restore the built-in layout.

\begin{tipbox}
Dans des conditions difficiles, une barre supérieure maigre avec widgets 3-4 clés (position, SOG, DPT, BTW)
est plus sûr que douze articles.
Ajouter des widgets seulement s'ils ajoutent une véritable conscience de la situation au passage actuel
phase.
\end{tipbox}

\chapter{La barre inférieure}
\label{ch:bottombar}

La barre inférieure est la bande d'action permanente en bas de chaque écran.
Its spatial stability is intentional: the \iconlg{applications}~Apps button is always
in the same place, the \iconlg{alarm-pob}~POB button is always within reach.
La mémoire musculaire, c'est la sécurité.

\screenshot{bottombar-overview}{Barre inférieure : bande d'applications ouvertes, icônes principales, état Signal~K}{}

\section{Régions à barre inférieure}

\rowcolors{2}{LtGray}{white}
\begin{longtable}{L{3.6cm}L{10.1cm}}
\toprule
\rowcolor{Navy}
\color{white}\sffamily\bfseries Region &
\color{white}\sffamily\bfseries Description \\
\midrule
Ouvrir la bande web-apps&
Une rangée de petites icônes d'application à gauche, une par application chargée.
Appuyez sur n'importe quelle icône pour sauter sur cette application instantanément.
  The active application is highlighted. \\
Icônes principales&
Le groupe central de sept icônes fixes.
  Spatially stable across all screens and comfort presets. \\
Zone d'état du signal ~K&
Côté droit. Indicateurs de santé compacts et REST montrant le nom du serveur,
dernière mise à jour timestamp, et état de connexion.
  Tapping opens a status detail drawer. \\
\bottomrule
\end{longtable}

\section{Référence principale de l'icône}

These are the exact icons defined in \code{BottomBar.ui}.

\rowcolors{2}{LtGray}{white}
\begin{longtable}{C{1.4cm}L{2.4cm}L{5.4cm}L{4.5cm}}
\toprule
\rowcolor{Navy}
\color{white}\sffamily\bfseries Icon &
\color{white}\sffamily\bfseries Label &
\color{white}\sffamily\bfseries Function &
\color{white}\sffamily\bfseries Notes \\
\midrule
\iconlg{applications}     & Apps      &
Ouvre le lanceur. Cache l'application actuelle et affiche la grille de tuiles.&
  Use as the safe escape hatch from any stuck or confusing screen. \\
\iconlg{alarm-pob}        & POB       &
Personne à bord. Un seul robinet déclenche l'urgence.&
  Always visible. Never disabled. \\
\iconlg{command-autopilot}& Autopilot &
Ouvre la barre de pilotage automatique.&
  Grey/inactive if no compatible autopilot plugin is detected. \\
\iconlg{anchor-iec}       & Anchor    &
Ouvre la barre Anchor Watch.&
  Grey/inactive if no anchor plugin is detected. \\
\iconlg{alarm-general}    & Alarms    &
Ouvre la barre des alarmes. En évidence quand une alarme est active.&
  Accent colour when active; draws the eye across all presets. \\
\iconlg{database}         & My~Data   &
  Opens the resource manager (Chapter~\ref{ch:mydata}). & \\
\iconlg{settings-iec}     & Settings  &
  Opens all configuration (Chapter~\ref{ch:settings}). & \\
\bottomrule
\end{longtable}

\section{Barres de recouvrement opérationnelles}

Tapping \iconlg{alarm-pob}~POB, \iconlg{command-autopilot}~Autopilot,
\iconlg{anchor-iec}~Anchor, or \iconlg{alarm-general}~Alarms replaces the normal Bottom
icônes de barre avec une barre opérationnelle dédiée.
La zone d'application continue de courir derrière la barre.
Tap \iconlg{close-google}~close (rightmost icon) to return to the normal Bottom Bar
sans changer d'état opérationnel.

\begin{warningbox}
Les barres opérationnelles ne sont pas modales.
Les données continuent à être mises à jour.
Le traceur de cartes et toutes les applications restent en direct derrière la barre opérationnelle.
\end{warningbox}

\chapter{Le lanceur et la gestion des applications}
\label{ch:launcher}

Le lanceur est la position sûre de la coque.
Utilisez-le pour revenir aux tuiles de l'application, basculer entre les tâches, laisser une application gelée, ou
vous réorienter après un flux de travail déroutant.
Chaque fois que l'écran actuel cesse de se sentir clair ou digne de confiance, appuyez sur
\iconlg{applications}~\key{Apps} to return here.

\screenshot{launcher-overview}{Écran d'accueil du lanceur: grille de tuiles, indicateur de page, Barre inférieure}{}

\section{Ouverture d'une demande}

\begin{enumerate}
  \item Tap \iconlabel{applications}{Apps} in the Bottom Bar.
  \item Navigate pages by swiping if needed.
  \item Tap the tile for the application you want.
  \item Wait for it to load in the Application Area.
\end{enumerate}

Si une application n'apparaît pas correctement, évitez les coups rapides répétés.
Attendez brièvement, essayez une fois si le shell offre l'option, puis retournez au lanceur
si l'application semble toujours mal.

\section{Changement d'applications ouvertes}

Utilisez la bande d'applications ouvertes dans la barre inférieure pour basculer entre les applications chargées
avec un seul robinet.
La bande montre une petite icône pour chaque application chargée; la
souligné.

\section{Ajout d'une application personnalisée}

\begin{enumerate}
  \item Open \ui{Settings~> Applications}.
  \item Tap \iconlabel{define-new-application-in-palette}{Create app}.
  \item Enter a display name and the full URL (\code{http://} or \code{https://}).
  \item Optionally set a custom icon path and zoom level.
  \item Tap \key{Save}.
\end{enumerate}

\section{Récupération d'une application en panne ou en panne}

\begin{enumerate}
  \item Wait briefly. Retry once if the shell offers a retry action.
  \item Tap \iconlabel{applications}{Apps} to return to the launcher.
  \item Check shell health and connection state.
  \item Reopen the app only after confirming the shell itself is healthy.
\end{enumerate}

\begin{tipbox}
Une demande rejetée ne signifie pas que FairWindSK a échoué.
Le shell, les widgets de données et les barres opérationnelles (POB, ancre, pilote automatique) continuent
travailler indépendamment de toute demande individuelle.
Une application dégradée n'est pas la même qu'un état dégradé du navire.
\end{tipbox}
